\documentclass[12pt,a4paper]{article}
\usepackage{xepersian}
\usepackage{graphicx}
\usepackage{hyperref}
\usepackage{listings}
\usepackage{color}
\usepackage{algorithm}
\usepackage{algpseudocode}
\usepackage{amsmath}
\usepackage{xskak}
\usepackage{chessboard}
\usepackage{minted}
\usepackage{tikz}
\usetikzlibrary{positioning, arrows.meta}
\usepackage{import}
\usepackage{natbib}   % or biblatex
\subimport{./Architecture_Graph/plotNeuralNet/layers/}{init}
\usetikzlibrary{positioning}

\def\ConvColor{rgb:yellow,5;red,2.5;white,5}
\def\ConvReluColor{rgb:yellow,5;red,5;white,5}
\def\PoolColor{rgb:red,1;black,0.3}
\def\DcnvColor{rgb:blue,5;green,2.5;white,5}
\def\SoftmaxColor{rgb:magenta,5;black,7}
\def\SumColor{rgb:blue,5;green,15}
\def\skipshift{6.5}
\def\poolsep{1}
\def\sep{3}
\def\cwidth{3}
\def\FcColor{rgb:blue,5;red,2.5;white,5}
\def\FcReluColor{rgb:blue,5;red,5;white,4}
\defpersianfont\nastaliq{IranNastaliq.ttf}
\settextfont[Scale = 1.0 ,
BoldFont = *Bd ,
ItalicFont = *It ,
BoldItalicFont = *BdIt ,
Extension = .ttf
]{XB Yas} 

\lstset{
	frame=single,
	breaklines=true,
	numbers=left,
	numberstyle=\tiny,
	basicstyle=\ttfamily\small,
	language=Python,
	xleftmargin=0pt,           % Remove left margin
	framexleftmargin=0pt,      % Remove frame margin
	linewidth=\textwidth,      % Set width to text width
	aboveskip=0pt,
	belowskip=0pt,
	resetmargins=true          % Important for Persian!
}

\title{گزارش پروژه موتور شطرنج با یادگیری عمیق و جستجوی درخت مونت کارلو}
\usepackage{xcolor}
\definecolor{codegreen}{rgb}{0,0.6,0}
\definecolor{codegray}{rgb}{0.5,0.5,0.5}
\definecolor{codepurple}{rgb}{0.58,0,0.82}
\begin{document}
	\maketitle
	\section{مقدمه}
	هدف از این پروژه، طراحی و پیاده‌سازی یک موتور شطرنج با استفاده از یادگیری عمیق و الگوریتم جستجوی درخت مونت کارلو (MCTS) است. این موتور قادر است با استفاده از یک شبکه عصبی کانولوشن، موقعیت‌های صفحه شطرنج را ارزیابی کرده و با ترکیب آن با جستجوی 
\lr{MCTS}
	،
	 سعی کند تا بهترین حرکت را انتخاب کند. داده‌های آموزش از پایگاه داده Lichess دریافت و پس از فیلتر کردن بر اساس ریتینگ بازیکنان، برای آموزش شبکه عصبی استفاده می‌شوند.
	
	\section{معماری شبکه عصبی}
\begin{figure}[h]
	\centering
	\includegraphics[width=0.8\textwidth]{my.pdf}
	\caption{معماری شبکه}
	\label{fig:architecture}
\end{figure}
	\subsection{ورودی شبکه}
	ورودی شبکه یک تنسور 19×8×8 است که هر لایه آن اطلاعات متفاوتی از صفحه شطرنج را نمایش می‌دهد:
	\begin{itemize}
		\item کانال‌های 0-5: موقعیت مهره‌های سفید (پیاده، اسب، فیل، رخ، وزیر، شاه)
		\item کانال‌های 6-11: موقعیت مهره‌های سیاه
		\item کانال‌های 12-15: حق قلعه رفتن
		\item کانال 16: خانه آن پاسان
		\item کانال 17: نوبت حرکت
		\item کانال 18: تکرار موقعیت
	\end{itemize}
	
	\subsection{لایه‌های شبکه }
	شبکه عصبی دارای 6 بلاک اولیه 
\lr{Residual}
هست، که منظور از 
\lr{Residual}
این است که ورودی هر بلاک، با خروجی اش جمع خواهد شد. این سبب می‌شود که مشکل محو شدگی گرادیان ها برطرف شود.
هر کدام از این بلاک ها شامل دو شبکه کانولوشنی 
\lr{Relu}
هستند، هر لایه کانولوشن از فیلترهای 3×3 با padding
برابر 1 استفاده می‌کند که باعث حفظ ابعاد 8×8 صفحه می‌شود. لایه‌های BatchNorm برای نرمال‌سازی استفاده شده‌اند.. نتایج این شبکه سپس به دو سر سیاست
\lr{ (Policy Head)}
و سر ارزش
\lr{(Value Head)}
داده می‌شود. که هدف سر سیاست، نگاه به صفحه شطرنج و تعیین احتمال هر حرکت در آن موقعیت است، و هدف سر ارزش، تعیین برتری بازیکن سفید یا سیاه است، که بین -1 و 1 قرار دارد. هر چه نزدیک تر به یک باشد، یعنی کاربر سفید برتری بیشتری دارد و هر چه نزدیک تر به -1 باشد یعنی کاربر سیاه برتری بیشتری دارد
. این معماری مشابه شبکه‌های استفاده شده در موتورهای معروفی مثل 
\lr{
 AlphaZero
\citep{alphazero}
}
 است اما با ابعاد کوچک‌تر برای اجرا روی سخت‌افزار معمولی و برخلاف 
\lr{AlphaZero}
از 
\lr{Supervised Learning}
استفاده می‌شود.
\\
کد معماری شبکه در 
\lr{pytorch}:
	\begin{latin}
	\begin{minted}[fontsize=\small, linenos]{python}
class ResidualBlock(nn.Module):
    def __init__(self, channels):
        super().__init__()
    
        self.conv1 = nn.Conv2d(channels, channels, kernel_size=3, padding=1)
        self.bn1 = nn.BatchNorm2d(channels)
        
        self.conv2 = nn.Conv2d(channels, channels, kernel_size=3, padding=1)
        self.bn2 = nn.BatchNorm2d(channels)
    
    def forward(self, x):
        identity = x

        out = self.conv1(x)
        out = self.bn1(out)
        out = F.relu(out)
        
        out = self.conv2(out)
        out = self.bn2(out)
        
        out = out + identity   # Skip connection
        out = F.relu(out)
        
        return out

class ChessNet(nn.Module):
    def __init__(self, num_blocks=6):
        super(ChessNet, self).__init__()
        
        # Initial feature extraction
        self.input_conv = nn.Sequential(
        nn.Conv2d(19, 128, 3, padding=1),
        nn.BatchNorm2d(128),
        nn.ReLU(inplace=True),
        )
        
        # Residual tower
        self.residual_layers = nn.Sequential(
        *[ResidualBlock(128) for _ in range(num_blocks)]
        )
        
        # Policy Head
        self.policy = nn.Sequential(
        nn.Conv2d(128, 32, 1),
        nn.BatchNorm2d(32),
        nn.ReLU(inplace=True),
        nn.Flatten(),
        nn.Linear(32 * 8 * 8, 4096)
        )
        
        # Value Head
        self.value = nn.Sequential(
        nn.Conv2d(128, 32, 1),
        nn.BatchNorm2d(32),
        nn.ReLU(inplace=True),
        nn.Flatten(),
        nn.Linear(32 * 8 * 8, 256),
        nn.ReLU(inplace=True),
        nn.Linear(256, 1),
        nn.Tanh()
        )
        
    def forward(self, x):
        x = self.input_conv(x)
        x = self.residual_layers(x)
    
        return self.policy(x), self.value(x)
)
\end{minted}
	\end{latin}
	
	\subsection{سر سیاست 
		\lr{(Policy Head)}}
	سر سیاست وظیفه تخمین احتمال  تمامی حرکات در هر موقعیت  شطرنج را بر عهده دارد:
	\\
	\begin{latin}
\begin{minted}[fontsize=\small, linenos]{python}
self.policy = nn.Sequential(
nn.Conv2d(128, 32, 1),
nn.BatchNorm2d(32),
nn.ReLU(inplace=True),
nn.Flatten(),
nn.Linear(32 * 64, 4096)
)
\end{minted}
	\end{latin}
	خروجی این سر یک بردار 4096 تایی است که هر عنصر نشان‌دهنده احتمال یک حرکت خاص است. تعداد 4096 از ضرب 64 (خانه مبدأ) در 64 (خانه مقصد) به دست آمده است. برخی از این حرکات ممکن است غیر قانونی باشند، به همین خاطر قبل از گرفتن 
    softmax
    ابتدا حرکات غیرقانونی را منهی یک عدد منفی خیلی بزرگ میکنیم تا مقدار softmax آنها صفر و در نتیجه تاثیری در احتمال حرکات مجاز نداشته باشند.
	
	\subsection{سر ارزش
		 \lr{(Value Head)}}
	سر ارزش موقعیت را ارزیابی کرده و یک عدد بین $-1$ تا $ 1$ برمی‌گرداند که نشان‌دهنده برتری سفید یا سیاه است:
	\\
	\begin{latin}
\begin{minted}[fontsize=\small, linenos]{python}
self.value = nn.Sequential(
nn.Conv2d(128, 32, 1),
nn.BatchNorm2d(32),
nn.ReLU(inplace=True),
nn.Flatten(),
nn.Linear(32 * 64, 256),
nn.ReLU(inplace=True),
nn.Linear(256, 1),
nn.Tanh()
)
\end{minted}
	\end{latin}
	\section{ارزیابی هیوریستیک}
	برای برچسب‌زنی داده‌های آموزش بدون نیاز به موتور های شطرنج مانند
	\lr{Stockfish}
	، یک تابع ارزیابی هیوریستیک استفاده شده است که عوامل مختلفی را در نظر می‌گیرد:
	\\
	\subsection{ارزش مهره‌ها}
	\begin{latin}
\begin{minted}[fontsize=\small, linenos]{python}
piece_values = {
    chess.PAWN: 100, chess.KNIGHT: 320, chess.BISHOP: 330,
    chess.ROOK: 500, chess.QUEEN: 900, chess.KING: 0
}
\end{minted}
	\end{latin}
	\subsection{جداول مکان مهره‌ها}
	برای هر مهره، یک جدول 8×8 از امتیازات موقعیتی تعریف شده است. به عنوان مثال، جدول پیاده‌ها به طور زیر است. که این یعنی هر چه پیاده به سطر آخر نزدیک تر باشد این بهتر است، چون رسیدن به سطر آخر به معنی ارتقا پیاده به مهره ای با ارزش بیشتر است. دلیل صفر بودن سطر آخر این است که وقتی پیاده به سطر آخر رسید، ارتقا خواهد یافت و دیگر پیاده نخواهد بود، پس هیچ وقت پیاده ای در سطر آخر وجود نخواهد داشت، پس مقادیر سطر آخر مهم نیست. برای بازیکن سیاه همین مقدار منفی خواهند شد و سطر هابرعکس می‌شوند.
	\\
	\begin{latin}
\begin{minted}[fontsize=\small, linenos]{python}
pawn_table = [
    [0,  0,  0,  0,  0,  0,  0,  0],
    [50, 50, 50, 50, 50, 50, 50, 50],
    [10, 10, 20, 30, 30, 20, 10, 10],
    [5,  5, 10, 25, 25, 10,  5,  5],
    [0,  0,  0, 20, 20,  0,  0,  0],
    [0,  0,  0,  0,  0,  0,  0,  0],
    [0,  0,  0,  0,  0,  0,  0,  0],
    [0,  0,  0,  0,  0,  0,  0,  0] 
    ]
\end{minted}
	\end{latin}
	\subsection{سایر عوامل}
	عوامل دیگری که در ارزیابی در نظر گرفته شده‌اند:
	\begin{itemize}		
		\item تحرک مهره‌ها (تعداد حرکات قانونی)
		\item کنترل مرکز صفحه
		\item ایمنی شاه و حق قلعه رفتن
		\item سپر پیاده‌ای جلوی شاه
	\end{itemize}
	
	نهایتاً امتیاز نهایی با تابع tanh نرمال‌سازی می‌شود:
	\[
	\text{value} = \tanh(\frac{\text{score}}{1000})
	\]
	
	\section{جستجوی درخت مونت کارلو (MCTS)}
	الگوریتم MCTS با ترکیب شبکه عصبی، جستجوی هوشمندانه‌ای برای یافتن بهترین حرکت انجام می‌دهد. این الگوریتم از 4 مرحله تشکیل شده است، که این مراحل تکرار خواهند شد(در کد تعداد پیش فرض تکرار 800 است).
	
	\subsection{مراحل MCTS}
	
	\subsubsection{انتخاب (Selection)}
	در این مرحله، با استفاده از فرمول PUTC، بهترین گره برای ادامه جستجو انتخاب می‌شود:
	\[
	\text{score} = Q(s,a) + c \times P(s,a) \times \frac{\sqrt{N(s)}}{1 + N(s,a)}
	\]
	
	که در آن:
	\begin{itemize}
		\item \(Q(s,a)\): برتری موقعیت شطرنج بعد از انجام حرکت 
        \lr{a}
        برای حرکت کننده.
		\item \(P(s,a)\): احتمال اولیه حرکت از شبکه عصبی که توسط سر سیاست شبکه تعیین می‌شود
		\item \(N(s)\): تعداد بازدیدهای گره والد
		\item \(N(s,a)\): تعداد دفعاتی که این حرکت انتخاب شده
		\item \(c\): ثابت اکتشاف (معمولاً \lr{1.4}) 
        که تعیین می‌کند موتور شطرنج چقدر مشتاق به اکتشاف حرکات جدید است.
	\end{itemize}
	
	\subsubsection{بسط (Expansion)}
	اگر گره انتخاب شده قبلاً بسط داده نشده باشد، تمام حرکات قانونی به عنوان فرزندان آن ایجاد می‌شوند. احتمال اولیه هر حرکت از خروجی Policy Head شبکه عصبی گرفته می‌شود.
	
	\subsubsection{ارزیابی (Evaluation)}
	اگر گره به حالت پایان بازی رسیده باشد، نتیجه واقعی بازی برگردانده می‌شود. در غیر این صورت، شبکه عصبی برای ارزیابی موقعیت فراخوانی می‌شود و خروجی سر ارزش به عنوان ارزش موقعیت در نظر گرفته می‌شود.
	
	\subsubsection{بازگشت (Backup)}
	ارزش به‌دست آمده در تمام گره‌های مسیر جستجو به‌روزرسانی می‌شود:
	\begin{itemize}
		\item تعداد بازدیدها (\(visits\)) یک واحد افزایش می‌یابد
		\item مجموع ارزش‌ها (\(value\_sum\)) با ارزش جدید جمع می‌شود
		\item برای گره‌های مربوط به حریف، علامت ارزش معکوس می‌شود
	\end{itemize}
	
	\subsection{پیاده‌سازی MCTS}
	کلاس‌های اصلی MCTS به صورت زیر پیاده‌سازی شده‌اند:
	\\
	\begin{latin}
\begin{minted}[fontsize=\small, linenos]{python}
class MCTSNode:
    def __init__(self, board, parent=None, move=None, prior=0):
        self.board = board
        self.parent = parent
        self.move = move
        self.prior = prior
        self.visits = 0
        self.value_sum = 0
        self.children = {}
        self.is_expanded = False

    def value(self):
        if self.visits == 0:
            return 0
        return self.value_sum / self.visits

    def ucb_score(self, exploration_constant=1.4):
        if self.visits == 0:
            return float('inf')
        exploration = exploration_constant * self.prior * \
            math.sqrt(self.parent.visits) / (1 + self.visits)
        return self.value() + exploration
\end{minted}
\end{latin}
	
	و کلاس اصلی MCTS:
	\\
	\begin{latin}
\begin{minted}[fontsize=\small, linenos]{python}
class MCTS:
    def __init__(self, model, num_simulations=800):
        self.model = model
        self.num_simulations = num_simulations
        self.device = torch.device("cuda" if torch.cuda.is_available() else "cpu")
        
    def search(self, board):
        root = MCTSNode(board.copy())
        for _ in range(self.num_simulations):
            node = root
            path = [node]
            # Selection
            while node.is_expanded and node.children:
                node = max(node.children.values(), key=lambda c: c.ucb_score())
                path.append(node)
            
            # Expansion
            if not node.board.is_game_over() and not node.is_expanded:
                self.expand(node)
                node.is_expanded = True
            
            # Evaluation
            if node.board.is_game_over():
                # Game over
                result = node.board.result()
                if result == "1-0":
                    value = 1
                elif result == "0-1":
                    value = -1
                else:
                    value = 0.0
            else:
                # Use neural network for evaluation
                value = self.evaluate(node.board)
            if node.board.turn == chess.WHITE:
                value = -value
            #Backpropagation
            for node in reversed(path):
                node.visits += 1
                node.value_sum += value
                value = -value  # Flip for next parent
                
        return root
\end{minted}
	\end{latin}
	
	\section{آماده‌سازی داده‌ها}
	داده‌های آموزش از پایگاه داده Lichess دریافت می‌شوند. این پایگاه داده شامل میلیون‌ها بازی شطرنج با فرمت PGN یا CVS است.
	
	\subsection{دریافت و فیلتر کردن فایل‌ها}
	فایل‌های PGN و CSV با فرمت فشرده zst از آدرس \lr{https://database.lichess.org/} دریافت می‌شوند. در این پروژه از سه فایل استفاده شده است:
	\begin{itemize}
		\item \lr{lichess\_db\_standard\_rated\_2015-01.pgn.zst} 
		\item \lr{lichess\_db\_standard\_rated\_2013-01.pgn.zst} 
        \item \lr{lichess\_db\_puzzle.csv.zst}
	\end{itemize}
	
	برای فیلتر کردن بر اساس ریتینگ، تابع \lr{safe\_elo} نوشته شده است. دو فایل اول شامل بازی های مختلف بین افراد مختلف است، فایل سوم شامل یکسری پازل شطرنج، همراه با راه حل آن هاست. 80 درصد این داده های برای train مدل و 20 درصد برای validation به کار رفته است.
	
	\subsection{استخراج موقعیت‌ها}
	از هر بازی، تعدادی موقعیت تصادفی (پیش‌فرض 15موقعیت) استخراج می‌شود. برای هر موقعیت، موارد زیر ذخیره می‌گردد:
	\begin{itemize}
		\item تنسور صفحه شطرنج (19×8×8)
		\item ارزش موقعیت از تابع ارزیابی هیوریستیک
		\item حرکتی که در آن موقعیت انجام شده که با عنوان هدف 
		\lr{policy}
		از آن استفاده می‌شود.
	\end{itemize}
	
	\section{آموزش مدل}
	فرآیند آموزش با استفاده از بهینه‌ساز AdamW و تابع Loss ترکیبی انجام می‌شود.
	
	\subsection{تابع Loss}
	تابع Loss از دو بخش تشکیل شده است:
	\[
	\mathcal{L} = \mathcal{L}_{\text{value}} + \mathcal{L}_{\text{policy}}
	\]
	
	که در آن:
	\begin{itemize}
		\item \(\mathcal{L}_{\text{value}} = \text{MSE}(v_{\text{pred}}, v_{\text{target}})\) خطای میانگین مربعات برای سر ارزش
		\item \(\mathcal{L}_{\text{policy}} = \text{CrossEntropy}(p_{\text{pred}}, p_{\text{target}})\) آنتروپی متقاطع برای سر سیاست
	\end{itemize}
	
	\subsection{جلوگیری از بیش‌برازش}
	برای جلوگیری از بیش‌برازش (Overfitting)، چندین روش استفاده شده است:
	
	\subsubsection{Weight Decay}
	در بهینه‌ساز AdamW از پارامتر weight\_decay با مقدار 1e-4 استفاده شده است. این کار یک جریمه به وزن‌های بزرگ اضافه می‌کند و باعث ساده‌تر شدن مدل می‌شود.
	
	\subsubsection{Early Stopping}
	 در هر epoch  بهترین مدل بر اساس کمترین خطای اعتبارسنجی انتخاب می‌گردد و ذخیره می‌شود:
	\\
	\begin{latin}
\begin{minted}[fontsize=\small, linenos]{python}
if avg_val < best_val_loss:
    best_val_loss = avg_val
    torch.save(self.model.state_dict(), 
        f"chess_model_BEST_val_{best_val_loss:.4f}.pth")
\end{minted}
	\end{latin}
	
	\section{اجرای برنامه}
ابتدا پکیج های مورد نیاز برناهه به طور زیر نصب می‌شوند:
\\

		\begin{latin}
		\begin{lstlisting}[language=Python]
python -m pip install -r requirements.txt
		\end{lstlisting}
	\end{latin}
سپس مدل را با دستور زیر می‌شود اجرا کرد:
\\
		\begin{latin}
	\begin{lstlisting}[language=Python]
python ./chess_engine_main.py
	\end{lstlisting}
\end{latin}

	برنامه دارای یک منوی تعاملی با گزینه‌های زیر است:
	\\
	\begin{latin}
		\begin{verbatim}
			Options:
			1. Train model
			2. Load existing model
			3. Play against engine
			4. Watch engine vs random
			5. Watch engine self-play
			6. Exit
		\end{verbatim}
	\end{latin}
	
	که اول گزینه دو برای لود کردن مدل تمرین داده شده و سپس می‌شود با مدل بازی کرد.
	برای تمرین دادن مدلی جدید باید ابتدا دستور:
	\\
			\begin{latin}
		\begin{lstlisting}[language=Python]
python ./data.py
		\end{lstlisting}
	\end{latin}

اجرا شود، این فایل داده های مورد نیاز برای تمرین مدل را آماده و در پوشه:
\lr{chess\_data}
ذخیره می کند.  سپس با اجرای فایل 
\lr{chess\_engine\_main.py}
و انتخاب گزینه 1 برای  تمرین مدل شروع می شود و مدل های یادگیری شده در فایل هایی با پسوند 
\lr{pth}
ذخیره می‌شوند که آن ها را میتوان  با استفاده از گزینه 2 لود کرد.
	\subsection{آماده‌سازی داده‌ها}
	برای آماده‌سازی داده‌ها، اسکریپت \lr{data.py} اجرا می‌شود.
	
	این اسکریپت فایل‌های zst را پردازش کرده و فایل‌های دسته‌ای با نام \lr{chess\_data\_batch\_*.pt} را در ایجاد میکند، که برای یادگیری استفاده می‌شوند.
	
	\subsection{آموزش مدل}
	پس از آماده‌سازی داده‌ها، با انتخاب گزینه 1 از منو و وارد کردن پارامترهای آموزش، فرآیند آموزش آغاز می‌شود. در طول آموزش، مدل‌ها با نام‌هایی مانند زیر ذخیره می‌شوند:
	\begin{itemize}
		\item \lr{chess\_model\_BEST\_val\_2.0604.pth} 
		\item \lr{chess\_model\_BEST\_val\_2.0529.pth} (بهترین مدل تا الان)
	\end{itemize}
	
	\subsection{بازی با موتور}
	با انتخاب گزینه‌های 3 تا 5 می‌توان با موتور بازی کرد یا آن را در مقابل حرکات تصادفی مشاهده نمود. موتور با انجام 800 شبیه‌سازی MCTS برای هر حرکت، بهترین حرکت را انتخاب می‌کند.
	\subsection{نتایج}
	این مدل با استفاده از
	\lr{ 11.5 }
	 میلیون موقعیت بازی شطرنج برای 
	\lr{14 epoch}
	 یادگیری داده شده. در هنگام بازی، تعداد شبیه سازی های MCTS برابر 800 قرار داده شده. در ادامه برخی بازی های این موتور در برابر حرکات شانسی نمایش داده شده (سفید موتور است و سیاه حرکات شانسی):
	 
\setLR
\begin{latin}
\newgame
\mainline{
	 1.e4 Nc6
	 2.Nf3 e6
	 3.d4 Nh6
	 4.Nc3 Ba3
	 5.bxa3 Qf6
	 6.Bd3 Qd8
	 7.O-O Ne5
	 8.Nxe5 c5
	 9.f4 g5
	 10.Qh5 b5
	 11.Qxf7+ Nxf7
	 12.Nb5 Rg8
	 13.d5 gxf4
	 14.Bxf4 exd5
	 15.exd5 Bb7
	 16.Nd6+ Ke7
	 17.Nf5+ Kf6
	 18.Nxf7 Rg4
	 19.Bg5+ Kg6
	 20.Nf5h6 Kh5
	 21.Bxd8 d6
	 22.Rf5 Kg6
	 23.Rg5#
}
\\
\chessboard


\textbf{white wins}
\end{latin}
\setRL
بازی دوم:
\setLR
\begin{latin}
\newgame
\mainline{
1.e4 g5
2.d4 a6
3.Nf3 c6
4.Bc4 Qb6
5.O-O Qc5
6.dxc5 e5
7.Bg5 f5
8.exf5 Ra7
9.Ne5 Ra8
10.Qh5
}
\\
\chessboard


\textbf{white wins}
\end{latin}
\setRL
بازی سوم:
\setLR
\begin{latin}
	\newgame
	\mainline{
1.e4 Nh6
2.Nf3 b5
3.Bxb5 g5
4.Nxg5 Ba6
5.Bxa6 Ng8
6.Bc4 c5
7.Bxf7# 
	}
	\\
	\chessboard
	
	
	\textbf{white wins}
\end{latin}
\setRL
در ادامه بازی ای داریم که موتور  در برابر خودش بازی کرده:
\setLR
\begin{latin}
	\newgame
	\mainline{
1.e4 Nf6
2.e5 Nd5
3.Nf3 Nc6
4.Bc4 Nb6
5.Bb3 d5
6.exd6 Qxd6
7.O-O Bg4
8.d4 Bxf3
9.Qxf3 Nxd4
10.Qxf7+ Kd8
11.Nc3 Nxb3
12.axb3 Kd7
13.Be3 Qc6
14.Rad1+ Kc8
15.Bxb6 axb6
16.Rfe1 e6
17.Ne4 Ra5
18.Rd3 Re5
19.Qf4 Rxe4
20.Rxe4 Qxc2
21.Rxe6 Qxd3
22.Re8+ Kd7
23.Re1 Bd6
24.Qf7+ Kc6
25.Qxg7 Rd8
26.Qf7 Rd7
27.Qe8 Qxb3
28.Qe4+ Kb5
29.Qxb7 Qxb2
30.Qc8 Kc6
31.Re8 Rf7
32.Qb8 Qxf2+
33.Kh1 Qf1#
}
	\\
	\chessboard
    \\
	\textbf{black wins}
\end{latin}

\bibliographystyle{plainnat}  % or another style
\bibliography{refs}  
\end{document}